\section{The Fourth-Moment Frontier}

Registered ten-cycles can be recognized without first projecting them
to quotient pentagons. Their indicator vectors carry an exact
fourth-order walk signature. This section identifies that signature,
determines the smallest nearby comparison class, and localizes the
difference to a concrete family of external relay paths.

\subsection{Cycle indicators and walk moments}

Let \(C\) be a simple ten-cycle in \(G_{60}\), and let
\[
x_C\in\mathbb R^{60}
\]
be its vertex-indicator vector:
\[
(x_C)_v
=
\begin{cases}
1, & v\in V(C),\\
0, & v\notin V(C).
\end{cases}
\]

Let \(A\) be the native adjacency matrix and
\[
L=4I-A
\]
the native graph Laplacian.

For \(k\geq0\), define the adjacency and Laplacian moments
\[
m_k^A(C)=x_C^{\mathsf T}A^k x_C,
\qquad
m_k^L(C)=x_C^{\mathsf T}L^k x_C.
\]

The adjacency moment counts length-\(k\) walks whose initial and final
vertices both lie on the cycle. The Laplacian moment is an equivalent
polynomial statistic because the graph is regular.

\subsection{Exact fourth-order recognition}

The complete census of \(1824\) simple ten-cycles was evaluated against
these moment sequences.

For every registered cycle,
\[
m_4^A(C)=880
\]
and
\[
m_4^L(C)=2160.
\]

No nonregistered ten-cycle has either value.

\begin{theorem}
For a simple ten-cycle \(C\subseteq G_{60}\), the following are
equivalent:
\[
C\text{ is registered},
\]
\[
x_C^{\mathsf T}A^4x_C=880,
\]
and
\[
x_C^{\mathsf T}L^4x_C=2160.
\]
\end{theorem}

Thus registered closure admits an exact blind recognition law at
fourth order. The test uses only the native graph, the cycle vertex
set, and the adjacency or Laplacian operator.

The full registered moment profiles through order eight are
\[
\bigl(m_k^L\bigr)_{k=0}^8
=
(10,20,80,400,2160,12000,67520,382720,2181120)
\]
and
\[
\bigl(m_k^A\bigr)_{k=0}^8
=
(10,20,80,240,880,3040,11200,40960,153600).
\]

\subsection{The third-moment frontier}

The moments through order three do not distinguish the registered
cycles uniquely.

Besides the twenty-four registered cycles, there are
\[
120
\]
nonregistered ten-cycles sharing the same lower-order adjacency
profile through \(k=3\).

We call these cycles the \emph{third-moment frontier}.

Thus the smallest comparison class relevant to fourth-order
recognition consists of
\[
24+120=144
\]
cycles.

\begin{proposition}
Registered cycles are not recognized by adjacency moments through
order three. At order four, all \(120\) frontier impostors separate
from all twenty-four registered cycles.
\end{proposition}

\subsection{Fourth-order walk anatomy}

For a registered cycle, the total fourth adjacency moment is
\[
m_4^A=880.
\]

The endpoint-distance decomposition is
\[
d_0:280,
\qquad
d_1:300,
\qquad
d_2:220,
\qquad
d_3:60,
\qquad
d_4:20,
\qquad
d_5:0,
\]
where \(d_j\) counts four-step walks beginning and ending on the cycle
whose native endpoints are at graph distance \(j\).

The boundary-crossing decomposition is
\[
0:160,
\qquad
2:560,
\qquad
4:160.
\]

Several individual fourth-order statistics are exact selectors for the
registered cycles, including
\[
\text{boundary-crossing count at }2=560,
\]
\[
\text{distance-4 endpoint count}=20,
\]
\[
\text{distance-5 endpoint count}=0,
\]
\[
\text{three-outside-vertex count}=80,
\]
and
\[
\text{five-distinct-quotient count}=100.
\]

These equalities are alternative descriptions of the same exact
fourth-order class.

\subsection{Antipodal exclusion}

The sharpest comparison with the third-moment frontier occurs in the
long-distance endpoint classes.

For every registered cycle,
\[
(d_4,d_5)=(20,0).
\]

For every frontier cycle,
\[
(d_4,d_5)=(28,8).
\]

Thus the registered cycles contain no four-step walk whose endpoints
are both on the cycle and lie at native graph distance five.

\begin{theorem}
Within the \(144\)-cycle third-moment frontier, a cycle is registered
if and only if it excludes four-step antipodal return:
\[
d_5=0.
\]
Every nonregistered frontier cycle has exactly eight such walks.
\end{theorem}

All frontier antipodal walks have the same coarse form:

\begin{enumerate}
    \item the initial vertex lies on the cycle;
    \item all three intermediate vertices lie outside the cycle;
    \item the final vertex lies on the cycle;
    \item the walk crosses the cycle boundary exactly twice;
    \item its quotient projection visits five chambers.
\end{enumerate}

The distinction is therefore not distributed arbitrarily across the
fourth-order walk census. It is concentrated in one external long
return family.

\subsection{Suppression of sixteen external returns}

The total adjacency moments are
\[
880
\]
for a registered cycle and
\[
896
\]
for a frontier cycle.

The difference is
\[
896-880=16.
\]

A finer comparison shows that the sixteen missing walks divide as
\[
8
\]
external distance-four returns and
\[
8
\]
external distance-five antipodal returns.

There is no compensating surplus in another fourth-order class.

\begin{theorem}
Relative to every third-moment frontier cycle, a registered cycle
suppresses exactly sixteen external four-step returns:
\[
8
\text{ at endpoint distance }4
\]
and
\[
8
\text{ at endpoint distance }5.
\]
All other fourth-order walk classes agree.
\end{theorem}

The word ``suppresses'' is important. The audit does not support a
redistribution claim. The registered count is obtained by deletion of
the sixteen external long returns, not by moving them into other
endpoint classes.

\subsection{The external bridge family}

For each frontier cycle, the sixteen additional walks are supported on
a six-vertex external configuration.

The configuration contains:

\begin{enumerate}
    \item four endpoint vertices lying on the cycle;
    \item two middle vertices outside the cycle;
    \item four endpoint--middle edges;
    \item no edge joining the two middle vertices;
    \item no edge joining two endpoints inside the relay structure.
\end{enumerate}

The induced relay graph is therefore not \(K_{2,4}\). It is
\[
P_3\sqcup P_3,
\]
the disjoint union of two paths of length two.

Each middle vertex joins exactly two cycle endpoints and contributes
eight oriented four-step returns. The two middle vertices together
contribute
\[
8+8=16.
\]

\begin{theorem}[Two-relay theorem]
Every nonregistered cycle in the third-moment frontier carries an
external relay graph isomorphic to
\[
P_3\sqcup P_3.
\]
The two relay components generate exactly the sixteen fourth-order
walks absent from a registered cycle.
\end{theorem}

The two middle vertices lie in one quotient chamber. Their native
relation is always \(b\) or \(ab\), never \(a\).

Across the \(120\) frontier cycles, the middle-relation profile is
\[
b:60,
\qquad
ab:60.
\]

This again reproduces the exchanged native pair
\[
b\longleftrightarrow ab.
\]

\subsection{Registered and frontier structures}

The fourth-order distinction can now be stated geometrically.

A frontier impostor carries two independent answering relays:
\[
P_3\sqcup P_3.
\]

A registered cycle carries neither relay.

Equivalently,
\[
\text{frontier cycle}
=
\text{registered lower-order profile}
+
\text{two external relays}.
\]

The registered class is characterized by the absence of those relay
paths.

\subsection{Recognition hierarchy}

The sequence of recognition results is:

\[
\text{quotient projection}
\quad\Longrightarrow\quad
\text{pentagon traversed twice},
\]
\[
\text{fourth moment}
\quad\Longrightarrow\quad
m_4^A=880,
\]
\[
\text{long-distance anatomy}
\quad\Longrightarrow\quad
(d_4,d_5)=(20,0),
\]
\[
\text{external geometry}
\quad\Longrightarrow\quad
\text{no }P_3\sqcup P_3\text{ relay}.
\]

These are equivalent recognitions of the same twenty-four-cycle class
within the relevant finite census.

\subsection{Boundary}

The moment calculation is a finite graph-dynamical signature. It does
not identify a physical energy, mass, force, or time variable.

The established statement is narrower and exact:

\begin{quote}
Registered closure is the unique ten-cycle class that excludes the
two external answering relays and their sixteen associated four-step
returns.
\end{quote}
